%% scientific_vita.tex — Wissenschaftlicher Lebenslauf / Scientific Vita
%% for dissertation submission by Christoph Gerhardt
%% Compile: pdflatex scientific_vita.tex (twice for hyperlinks)
\documentclass[11pt,a4paper]{article}

% --- Font & Encoding ----------------------------------------------------------
\usepackage[utf8]{inputenc}
\usepackage[T1]{fontenc}
\usepackage{mathpazo}                   % Palatino — matching thesis font
\usepackage{microtype}

% --- Page Geometry (matching thesis) ------------------------------------------
\usepackage[a4paper, margin=3cm, top=3.8cm]{geometry}

% --- Spacing ------------------------------------------------------------------
\usepackage{setspace}
\onehalfspacing

% --- Packages -----------------------------------------------------------------
\usepackage[inline]{enumitem}
\usepackage{xcolor}
\usepackage{tabularx}
\usepackage{booktabs}
\usepackage{tabularx}
\usepackage{booktabs}
\usepackage[hyphens]{url}
\usepackage[
  colorlinks=true,
  linkcolor=tuilmenaugruen,
  citecolor=tuilmenaugruen,
  urlcolor=tuilmenaublau,
  bookmarks=true
]{hyperref}

% --- TU Ilmenau Corporate Colors ---------------------------------------------
\definecolor{tuilmenauorange}{HTML}{FF7900}
\definecolor{tuilmenaublau}{HTML}{003359}
\definecolor{tuilmenaugruen}{HTML}{00747A}

% --- Title Formatting ---------------------------------------------------------
\usepackage{titlesec}
\titleformat{\section}
  {\Large\bfseries\color{tuilmenaugruen}}
  {}
  {0em}
  {}
  [\vspace{-0.5em}\textcolor{tuilmenaugruen}{\rule{\textwidth}{0.8pt}}]

\titleformat{\subsection}
  {\large\bfseries\color{tuilmenaugruen}}
  {}
  {0em}
  {}

% --- CV entry commands --------------------------------------------------------
\newcolumntype{L}{>{\raggedright\arraybackslash}X}
\newcommand{\cvline}[2]{%
  \par\noindent
  \begin{minipage}[t]{3.5cm}\raggedright\textbf{#1}\end{minipage}%
  \hfill
  \begin{minipage}[t]{\dimexpr\textwidth-3.8cm\relax}#2\end{minipage}%
  \par\vspace{4pt}%
}
\newcommand{\cventryFull}[5]{%
  \par\noindent
  \begin{minipage}[t]{3.5cm}\raggedright\textbf{#1}\end{minipage}%
  \hfill
  \begin{minipage}[t]{\dimexpr\textwidth-3.8cm\relax}%
    \textbf{#2}\par #3\par\textit{#4}%
    \ifx&#5&\else\par #5\fi
  \end{minipage}%
  \par\vspace{8pt}%
}

% ==============================================================================
\begin{document}

% --- Header -------------------------------------------------------------------
\begin{center}
  {\LARGE\bfseries\color{tuilmenaugruen} Scientific Vita}\\[0.3em]
  {\large\color{tuilmenaugruen} (Wissenschaftlicher Lebenslauf)}\\[0.8em]
  {\large Christoph Gerhardt}\\[0.3em]
  {\normalsize born July 17, 1990 in Ilmenau, Germany}\\[1.5em]
\end{center}

% ==============================================================================
\section*{Education}
% ==============================================================================

\cventryFull{Oct 2019 -- May 2026\newline(expected)}
  {Dr.-Ing.\@ (PhD) in Computer Science}
  {Technische Universit\"at Ilmenau, Faculty of Computer Science and Automation}
  {Dissertation: ``Comprehensive Outdoor Scene Understanding Through Multi-Scale
   Neural Networks: From 360\textdegree{} Object Detection to Panoramic
   Segmentation and Weather-Robust Augmentation''}
  {Supervisor: Prof.\@ Dr.\@ rer.\@ nat.\@ Wolfgang Broll\newline
   Research Group: Virtual Worlds and Digital Games (VWDG)}

\cventryFull{Oct 2014 -- Sep 2017}
  {M.Sc.\@ Computer Engineering (Ingenieurinformatik)}
  {Technische Universit\"at Ilmenau --- Final Grade: 1.5 (very good)}
  {Specialisation: Multimedia Information and Communication Systems.\newline
   Thesis: ``Selektive Verschl\"usselung von Gesichtern in Videodaten
   unter Verwendung des H.264-Standards''}
  {}

\cventryFull{Oct 2010 -- Sep 2014}
  {B.Sc.\@ Computer Engineering (Ingenieurinformatik)}
  {Technische Universit\"at Ilmenau --- Final Grade: 1.8 (good)}
  {Thesis: ``Entwicklung eines Verfahrens zur Verbesserung von
   Mikrofonklassifizierung basierend auf der Analyse von Umwelteinfl\"ussen''}
  {}

\cventryFull{Jun 2009}
  {Abitur}
  {Gymnasium ``Am Lindenberg'', Ilmenau --- Grade: 1.5}
  {}
  {}

% ==============================================================================
\section*{Professional Experience}
% ==============================================================================

\cventryFull{Feb 2019 -- present}
  {Research Associate (Wissenschaftlicher Mitarbeiter)}
  {Technische Universit\"at Ilmenau, Fachgebiet Virtuelle Welten und Digitale Spiele}
  {Doctoral research on outdoor scene understanding with neural networks;
   development of datasets, architectures, and benchmarks; teaching; student supervision}
  {}

\cventryFull{Feb 2017 -- Jan 2019}
  {Software Developer}
  {Fraunhofer Institute for Digital Media Technology (IDMT), Ilmenau}
  {Group: Virtual Acoustics --- C++ and Python development for spatial audio rendering}
  {}

\cventryFull{Jul 2016 -- Jan 2017}
  {Student Research Assistant (Hilfswiss.\@ Mitarbeiter)}
  {Fraunhofer IDMT, Ilmenau}
  {Group: Media Distribution and Security --- Face detection and selective encryption in H.264 video}
  {}

% ==============================================================================
\section*{Research Activities}
% ==============================================================================

\subsection*{Dissertation Research (2019--2026)}

\noindent
My doctoral research addresses the challenge of \textbf{comprehensive outdoor scene
understanding} through a progression of increasingly complex perception tasks:

\begin{enumerate}[nosep, leftmargin=1.5em, itemsep=4pt]
  \item \textbf{360\textdegree{} Traffic Sign Detection} --- Adapted object detection networks
    for equirectangular panoramic images; created a dataset of 2\,500 images with
    81 traffic sign classes; deployed in NOVASIB GmbH's TT-SIB INFOSYS\textregistered{} road
    inventory system (IEEE ITSC 2020).

  \item \textbf{Universal Outdoor Scene Segmentation (OUTSIDE)} --- Designed a multi-scale
    semantic segmentation architecture for unconstrained outdoor scenes; curated the
    OUTSIDE15k dataset (15\,000 images, 24 classes) (IEEE MMSP 2021).

  \item \textbf{Sky and Cloud Segmentation (SkyCloud)} --- Developed SkyCloudNet for
    fine-grained sky and cloud segmentation from perspective images; released dataset
    and pre-trained models (IEEE ICIVC 2023).

  \item \textbf{Panoramic Sky Segmentation (SkyCloud360)} --- Extended sky/cloud
    segmentation to equirectangular 360\textdegree{} imagery; created a 600-image panoramic
    dataset with addressing of projection distortions (IEEE ICIVC 2025).

  \item \textbf{Weather-Robust Augmentation (AWARE)} --- Developed a three-pipeline
    benchmark framework (SWIFT, PRISM, PROVE) evaluating 25 augmentation strategies
    for weather robustness; generated the AWACS dataset (71\,400 images)
    (IEEE Access, under review 2025).
\end{enumerate}

\subsection*{Collaborative Research}

\noindent
In addition to the core dissertation work, I contributed to collaborative research
within the VWDG group on:

\begin{itemize}[nosep, leftmargin=1.5em, itemsep=3pt]
  \item \textbf{Avatar and Agent Realism} --- Behavioural realism of animated avatars in
    AR and VR environments, including single-point IK animation and full-body
    motion capture comparisons (IEEE VR 2024, Int.\@ Conf.\@ Cyberworlds 2024,
    The Visual Computer 2025).
  \item \textbf{AR Communication Systems} --- User studies with older adults using
    wearable AR communication systems (ACM IMX 2024, ITAP/HCII 2025).
  \item \textbf{Gaze and Audiovisual Coherence} --- Analysis of gaze behaviour and
    effects of visual realism and audio presentation in AR storytelling
    (ACM ETRA 2024, IEEE VR 2024).
  \item \textbf{Systematic Reviews} --- Co-authored a systematic review on avatar and
    agent visualisation in HMD-based AR \& VR (IEEE TVCG 2023).
  \item \textbf{Telepresence Quality} --- Impact of partner representation
    (robot, avatar, human) on communication experience (ACM CSCW 2025).
\end{itemize}

\subsection*{Open-Source Contributions \& Datasets}

\begin{itemize}[nosep, leftmargin=1.5em, itemsep=3pt]
  \item \textbf{AWARE Framework}: SWIFT + PRISM + PROVE pipelines
    (\url{https://carhartt21.github.io/AWARE})
  \item \textbf{AWACS Dataset}: 71\,400 weather-augmented images
  \item \textbf{SkyCloud360}: 600 equirectangular images
    (\url{https://github.com/carhartt21/SkyCloud360})
  \item \textbf{SkyCloudNet}: Sky/cloud segmentation models
    (\url{https://github.com/carhartt21/SkyCloudNet})
  \item \textbf{OUTSIDE15k}: 15\,000 images, 24 classes
    (\url{https://github.com/carhartt21/outsideNet})
  \item \textbf{360\textdegree{} Detection}: 2\,500 panoramic images, 81 sign classes
    (\url{https://github.com/carhartt21/360-Object-Detection})
\end{itemize}

% ==============================================================================
\section*{Conference Presentations}
% ==============================================================================

\begin{enumerate}[nosep, leftmargin=1.5em, itemsep=4pt]
  \item \textbf{Oral Presentation}: ``SkyCloud360: Sky and Cloud Segmentation in
    Equirectangular Images,'' 10th Int.\@ Conf.\@ Image, Vision and Computing (ICIVC),
    Chengdu, China, June 2025.

  \item \textbf{Oral Presentation}: ``SkyCloud: Neural Network-Based Sky and Cloud
    Segmentation from Natural Images,'' 8th Int.\@ Conf.\@ Image, Vision and Computing
    (ICIVC), Dalian, China, August 2023.

  \item \textbf{Oral Presentation}: ``OUTSIDE: Multi-Scale Semantic Segmentation of
    Universal Outdoor Scenes,'' IEEE 23rd Int.\@ Workshop Multimedia Signal Processing
    (MMSP), Tampere, Finland, October 2021.

  \item \textbf{Oral Presentation}: ``Neural Network-Based Traffic Sign Recognition in
    360\textdegree{} Images for Semi-Automatic Road Maintenance Inventory,'' IEEE 23rd Int.\@
    Conf.\@ Intelligent Transportation Systems (ITSC), virtual, September 2020.
\end{enumerate}

% ==============================================================================
\section*{Teaching Activities}
% ==============================================================================

\subsection*{Courses (Teaching Assistant)}

\noindent\begin{tabularx}{\textwidth}{@{}p{3.5cm}L@{}}
\textbf{WS 2019--2025} & \textbf{Virtual and Augmented Reality}\\
 & Technische Universit\"at Ilmenau\\
 & Practical sessions on VR/AR fundamentals, 3D rendering, tracking systems,
   and immersive interaction design. Student project mentoring and grading.\\[6pt]
\textbf{SS 2020--2025} & \textbf{Virtuelle Welten}\\
 & Technische Universit\"at Ilmenau\\
 & Practical sessions on virtual worlds, real-time rendering, game engine
   development, and 3D modelling. Semester project supervision.\\[6pt]
\textbf{SS 2026} & \textbf{Virtual and Augmented Reality}\\
 & Technische Universit\"at Ilmenau\\[6pt]
\textbf{WS 2025/26} & \textbf{Game Development}\\
 & Technische Universit\"at Ilmenau\\
\end{tabularx}

\subsection*{Seminar Supervision}

\noindent
Supervision of seminars on computer vision, deep learning, semantic segmentation,
object detection, weather-robust perception, and neural network architectures
(ongoing, Technische Universit\"at Ilmenau).

\subsection*{Master Thesis Supervision}

\begin{itemize}[nosep, leftmargin=1.5em, itemsep=3pt]
  \item \textbf{Zamam Mirza}: ``Synthetic Data Generation for Semantic
    Segmentation and Object Detection Using Generative Models''
  \item \textbf{Syed Ajaz Zaidi}: ``Enhancing Digital Reading with
    Gaze-Triggered Audio Cues and Audiobook Listening with Augmented Sound Effects''
\end{itemize}

\subsection*{Research Project Supervision}

\cvline{SS 2024}{%
  Generative Models for Data Augmentation}

\cvline{WS 2024/25}{%
  Mimic Synthesis for 3D Avatars using Gaze and Audio;
  Automatic Non-Reference Quality Assessment for Artificially Generated Images;
  Cross-Domain Image-to-Image Translation for Low-Light Image Enhancement;
  3D Avatar Generation from Monocular Video;
  Guided Image Generation Using Stable Diffusion Models;
  Realistic Hair Generation for 3D Avatars using Deep Learning}

\cvline{SS 2025}{%
  Data Augmentation in Adverse Outdoor Conditions using Simulator Data;
  Utilizing Simulator-Generated Data for Enhanced Model Training;
  Non-Reference Quality Assessment for Synthetic Images}

\cvline{WS 2025/26}{%
  Non-Reference Quality Assessment for Generated Images;
  Large Tutoring Models;
  Egocentric Video Pretraining for XR Assistants;
  LLM-Driven Facial Animations for Avatars}

\cvline{SS 2026}{%
  Image-to-Image Translation in Equirectangular/360\textdegree{} Images}

% ==============================================================================
\section*{Funding}
% ==============================================================================

\cvline{2019 -- 2024}{Research funded by the \textbf{Free State of Thuringia}
  (Grant No.\@ 2017\,FE\,9147), co-financed by the European Union under the
  European Regional Development Fund (ERDF). Project focus: intelligent road
  infrastructure analysis using neural networks and 360\textdegree{} imagery.}

% ==============================================================================
\section*{Industry Collaboration}
% ==============================================================================

\cvline{2019 -- 2022}{Collaboration with \textbf{NOVASIB GmbH} on the integration of
  neural-network-based traffic sign detection into the TT-SIB INFOSYS\textregistered{}
  road maintenance inventory management system, enabling semi-automatic
  road infrastructure assessment from 360\textdegree{} panoramic imagery.}

% ==============================================================================
\section*{Technical Skills}
% ==============================================================================

\cvline{Programming}{Python, C/C++/C\#, Java, MATLAB/Octave, SQL, Bash, \LaTeX}
\cvline{ML \& DL}{PyTorch, TensorFlow/Keras, scikit-learn, OpenCV, ONNX, Albumentations, timm}
\cvline{Infrastructure}{Git, Docker, Linux, Slurm HPC, Weights~\&~Biases, Jupyter, CI/CD}
\cvline{Domains}{Semantic Segmentation, Object Detection, Data Augmentation, 360\textdegree{} Vision, VR/AR}

% ==============================================================================
\section*{Languages}
% ==============================================================================

\cvline{German}{Native speaker}
\cvline{English}{Fluent (C1+) --- academic writing, conference presentations}
\cvline{French}{Basic knowledge (A2)}
\cvline{Serbian}{Basic knowledge (A2)}

% ==============================================================================
\section*{Certifications}
% ==============================================================================

\cvline{2023}{C-Lizenz (DFB Football Coaching License)}
\cvline{2023}{EU Remote Pilot Certificate (A1/A3) --- Drone operation, EU Reg.\@ 2019/947}

% ==============================================================================
\section*{Reviewer Activities}
% ==============================================================================

\subsection*{Journal Reviewing}

\begin{itemize}[nosep, leftmargin=1.5em, itemsep=3pt]
  \item \textbf{IEEE Transactions on Visualization and Computer Graphics (TVCG)}\\
    ``Towards Privacy-preserving Photorealistic Self-avatars in Mixed Reality''
  \item \textbf{IEEE Access}\\
    ``Embodiment and Presence in Virtual Reality: The Role of Self-Similarity
     in High-Fidelity 3D Avatars''
  \item \textbf{IEEE Open Journal of Signal Processing}\\
    ``Unsupervised Angularly Consistent 4D Light Field Segmentation using
     Hyperpixels and a Graph Neural Network''
\end{itemize}

\subsection*{Conference Reviewing}

\begin{itemize}[nosep, leftmargin=1.5em, itemsep=3pt]
  \item \textbf{IEEE ISMAR 2025}\\
    ``Taking Language Embedded 3D Gaussian Splatting into the Wild''
  \item \textbf{IEEE ISMAR 2024}\\
    ``Exploring Device-Oriented Video Encryption for Hierarchical
     Privacy Protection in AR Content Sharing''\\
    ``HOIMotion: Forecasting Human Motion During Human-Object Interactions
     Using Egocentric 3D Object Bounding Boxes''\\
    ``Comparative Analysis of AR-Based and Traditional Approaches in
     English Pronunciation Teaching: Efficacy, Engagement, and Implementation''
\end{itemize}

% ==============================================================================

\vfill
\begin{flushright}
  \small Ilmenau, \today
\end{flushright}

\end{document}
